\documentclass[11pt]{article}
\usepackage[a4paper,margin=25mm]{geometry}
\usepackage{amsmath,amssymb,hyperref}
\hypersetup{colorlinks=true,urlcolor=blue}
\title{Strict log-convexity and the normalized Gamma characterization\\A proof of Conjecture 00000003282}
\author{ziangni-sys}\date{10 October 2026}
\begin{document}\maketitle
\begin{abstract}
On the positive real axis, the Gamma function is strictly log-convex and is the unique positive log-convex function satisfying the factorial recurrence and the normalization at one. More generally, every positive log-convex recurrence solution is exactly $f(1)\Gamma$. We give the complete classical characterization, including an elementary strictness argument from the recurrence, and formalize it for Mathlib's actual Gamma function.
\end{abstract}

\section{Precise characterization}
Both versions of the source invoke the Bohr--Mollerup characterization and a uniqueness constant equal to one. Its standard, precise meaning is the following theorem on $(0,\infty)$.

\medskip\noindent\textbf{Theorem.} The function $\Gamma$ is positive, satisfies
\[
\Gamma(1)=1,\qquad \Gamma(x+1)=x\Gamma(x)\quad(x>0),
\]
and $\log\Gamma$ is strictly convex. If $f:(0,\infty)\to(0,\infty)$ satisfies the same recurrence and has convex logarithm, then
\[
f(x)=f(1)\Gamma(x)\quad(x>0).
\]
Consequently $f(1)=1$ makes $f=\Gamma$, with multiplicative constant exactly one. Equivalently, positivity, normalized recurrence, and strict log-convexity characterize $\Gamma$.

All these characterizing hypotheses are essential to the formulation: strict log-convexity alone is not a uniqueness condition. The theorem is a positive-real characterization, the usual meaning of Bohr--Mollerup; Gamma's meromorphic continuation is not log-convex on its complex domain. The factorial values follow immediately by induction from the normalized recurrence.

\section{Existence and strictness}
For $x>0$, Euler's integral is
\[
\Gamma(x)=\int_0^\infty t^{x-1}e^{-t}\,dt.
\]
It converges, is positive, and gives $\Gamma(1)=1$. Integration by parts gives the recurrence; the boundary terms vanish because $t^xe^{-t}\to0$ at infinity and $t^x\to0$ at zero.

For $x,y>0$ and $a,b>0$ with $a+b=1$, H\"older's inequality gives
\begin{align*}
\Gamma(ax+by)
 &=\int_0^\infty(t^{x-1}e^{-t})^a(t^{y-1}e^{-t})^b\,dt\\
 &\le \Gamma(x)^a\Gamma(y)^b.
\end{align*}
Taking logarithms proves convexity of $L(x)=\log\Gamma(x)$.

Here strictness follows without analyzing equality in H\"older. The recurrence says
\[
L(x)=L(x+1)-\log x.
\]
The first function on the right is convex, as a translate of $L$. The second is strictly convex on $(0,\infty)$, since $(-\log x)''=1/x^2>0$. Their sum is strictly convex. Explicitly, for $x\ne y$ and $a,b>0$, $a+b=1$, convexity of $L$ and strict concavity of the logarithm give
\begin{align*}
L(ax+by)
 &=L(a(x+1)+b(y+1))-\log(ax+by)\\
 &<aL(x+1)+bL(y+1)-a\log x-b\log y\\
 &=aL(x)+bL(y).
\end{align*}
This proves the strict log-convexity asserted in the conjecture.

\section{Uniqueness, including the constant}
Suppose $f>0$ has convex logarithm and satisfies the recurrence. Put $c=f(1)>0$. Induction yields $f(n)=c(n-1)!$ for positive integers $n$.

Fix $0<x\le1$ and an integer $n\ge2$. The increasing secant slopes of the convex function $\log f$ give
\[
\log f(n)+x\log(n-1)
 \le\log f(n+x)
 \le\log f(n)+x\log n.
\]
For the upper bound use convexity between $n$ and $n+1$; for the lower bound compare the slope from $n-1$ to $n$ with that from $n$ to $n+x$. Repeated recurrence also gives
\[
f(n+x)=f(x)\prod_{k=0}^{n-1}(x+k).
\]
Therefore, with
\[
A_n(x)=\frac{(n-1)!}{\prod_{k=0}^{n-1}(x+k)},
\]
we have the two-sided estimate
\[
cA_n(x)(n-1)^x\le f(x)\le cA_n(x)n^x.
\]
The already established properties of $\Gamma$ give the same estimates with $c=1$. Dividing positive quantities consequently yields
\[
\left(\frac{n-1}{n}\right)^x
\le\frac{f(x)}{c\Gamma(x)}
\le\left(\frac{n}{n-1}\right)^x.
\]
Both endpoints tend to one, so $f(x)=c\Gamma(x)$ for $0<x\le1$. For an arbitrary $y>0$, write $y=x+m$ with $0<x\le1$ and a nonnegative integer $m$. Repeated recurrence for both functions extends this identity to $y$. This proves the complete classification.

If $f(1)=1$, then $c=1$ and $f=\Gamma$. Conversely $\Gamma$ has every required property. A multiple $c\Gamma$ can have value one at one only when $c=1$, because $\Gamma(1)=1$. This is the precise normalization and uniqueness constant in the source.

\section{Formal proof and reproduction}
The Lean project uses Mathlib's actual \texttt{Real.Gamma}. It proves strict log-convexity from translated convexity and strict concavity of \texttt{Real.log}, verifies all characterizing properties, proves normalized uniqueness, and normalizes an arbitrary positive recurrence solution by division by $f(1)$ to prove the full multiplicative classification. Its final biconditional says that the full characterization holds exactly when the function agrees with Gamma at every positive real argument.

The formal proof uses Mathlib's established Gamma integral, recurrence, log-convexity, and Bohr--Mollerup theorem. That theorem is a proved library result, not an axiom or a hypothesis asserting uniqueness. The report above supplies the classical complete uniqueness argument. No claim is made about arbitrary values of the totalized Lean functions outside the positive real axis.

Run \texttt{lake build} in \texttt{lean/}, using Lean 4.19.0. The public Mathlib dependency is pinned to
\begin{center}\texttt{c44e0c8ee63ca166450922a373c7409c5d26b00b}.\end{center}
The build prints axiom audits for strictness, the Gamma properties, normalized uniqueness, the rescaling classification, the constant-one theorem, and the final characterization. No proof placeholders, custom axioms, or native evaluation shortcuts are used.
\end{document}
